% Options for packages loaded elsewhere
\PassOptionsToPackage{unicode}{hyperref}
\PassOptionsToPackage{hyphens}{url}
\documentclass[
  ignorenonframetext,
]{beamer}
\newif\ifbibliography
\usepackage{pgfpages}
\setbeamertemplate{caption}[numbered]
\setbeamertemplate{caption label separator}{: }
\setbeamercolor{caption name}{fg=normal text.fg}
\beamertemplatenavigationsymbolsempty
% remove section numbering
\setbeamertemplate{part page}{
  \centering
  \begin{beamercolorbox}[sep=16pt,center]{part title}
    \usebeamerfont{part title}\insertpart\par
  \end{beamercolorbox}
}
\setbeamertemplate{section page}{
  \centering
  \begin{beamercolorbox}[sep=12pt,center]{section title}
    \usebeamerfont{section title}\insertsection\par
  \end{beamercolorbox}
}
\setbeamertemplate{subsection page}{
  \centering
  \begin{beamercolorbox}[sep=8pt,center]{subsection title}
    \usebeamerfont{subsection title}\insertsubsection\par
  \end{beamercolorbox}
}
% Prevent slide breaks in the middle of a paragraph
\widowpenalties 1 10000
\raggedbottom
\AtBeginPart{
  \frame{\partpage}
}
\AtBeginSection{
  \ifbibliography
  \else
    \frame{\sectionpage}
  \fi
}
\AtBeginSubsection{
  \frame{\subsectionpage}
}
\usepackage{iftex}
\ifPDFTeX
  \usepackage[T1]{fontenc}
  \usepackage[utf8]{inputenc}
  \usepackage{textcomp} % provide euro and other symbols
\else % if luatex or xetex
  \usepackage{unicode-math} % this also loads fontspec
  \defaultfontfeatures{Scale=MatchLowercase}
  \defaultfontfeatures[\rmfamily]{Ligatures=TeX,Scale=1}
\fi
\usepackage{lmodern}
\ifPDFTeX\else
  % xetex/luatex font selection
\fi
% Use upquote if available, for straight quotes in verbatim environments
\IfFileExists{upquote.sty}{\usepackage{upquote}}{}
\IfFileExists{microtype.sty}{% use microtype if available
  \usepackage[]{microtype}
  \UseMicrotypeSet[protrusion]{basicmath} % disable protrusion for tt fonts
}{}
\makeatletter
\@ifundefined{KOMAClassName}{% if non-KOMA class
  \IfFileExists{parskip.sty}{%
    \usepackage{parskip}
  }{% else
    \setlength{\parindent}{0pt}
    \setlength{\parskip}{6pt plus 2pt minus 1pt}}
}{% if KOMA class
  \KOMAoptions{parskip=half}}
\makeatother
\usepackage{graphicx}
\makeatletter
\newsavebox\pandoc@box
\newcommand*\pandocbounded[1]{% scales image to fit in text height/width
  \sbox\pandoc@box{#1}%
  \Gscale@div\@tempa{\textheight}{\dimexpr\ht\pandoc@box+\dp\pandoc@box\relax}%
  \Gscale@div\@tempb{\linewidth}{\wd\pandoc@box}%
  \ifdim\@tempb\p@<\@tempa\p@\let\@tempa\@tempb\fi% select the smaller of both
  \ifdim\@tempa\p@<\p@\scalebox{\@tempa}{\usebox\pandoc@box}%
  \else\usebox{\pandoc@box}%
  \fi%
}
% Set default figure placement to htbp
\def\fps@figure{htbp}
\makeatother
\setlength{\emergencystretch}{3em} % prevent overfull lines
\providecommand{\tightlist}{%
  \setlength{\itemsep}{0pt}\setlength{\parskip}{0pt}}
% Slide layout defaults for warning-clean scientific decks.
\usepackage{etoolbox}
\IfFileExists{xurl.sty}{\usepackage{xurl}}{}
\IfFileExists{seqsplit.sty}{\usepackage{seqsplit}}{\newcommand{\seqsplit}[1]{#1}}
\protected\def\breakseq#1{\seqsplit{#1}}
\protected\def\breaktt#1{\begingroup\ttfamily\seqsplit{#1}\endgroup}
\setlength{\emergencystretch}{6em}
\tolerance=5000
\hbadness=10000
\hfuzz=1pt
\setlength{\tabcolsep}{2pt}
\AtBeginEnvironment{longtable}{\tiny\renewcommand{\arraystretch}{0.86}\setlength{\tabcolsep}{1pt}}
\AtBeginEnvironment{tabular}{\tiny\renewcommand{\arraystretch}{0.86}\setlength{\tabcolsep}{1pt}}

% Natbib and cross-reference fallbacks — slides don't load natbib
% or cleveref, but combined-PDF manuscript prose may emit these
% commands. The fallback renders citations as a bracketed key list
% and cross-references as detokenized labels so slides don't fail on
% undefined control sequences or raw underscores. \providecommand is
% a no-op if packages load later.
\providecommand{\citep}[1]{[#1]}
\providecommand{\citet}[1]{#1}
\providecommand{\citealp}[1]{#1}
\providecommand{\citeauthor}[1]{#1}
\providecommand{\citeyear}[1]{#1}
\providecommand{\cref}[1]{\texttt{\detokenize{#1}}}
\providecommand{\Cref}[1]{\texttt{\detokenize{#1}}}

\newtheorem{proposition}{Proposition}
\newtheorem{hypothesis}{Hypothesis}
\usepackage{bookmark}
\IfFileExists{xurl.sty}{\usepackage{xurl}}{} % add URL line breaks if available
\urlstyle{same}
% fallback for those not using the hyperref driver hyperxmp:
\makeatletter
\@ifundefined{xmpquote}{\newcommand{\xmpquote}[1]{#1}}{}
\makeatother
\hypersetup{
  hidelinks,
  pdfcreator={LaTeX via pandoc}}

\author{\texorpdfstring{}{}}
\date{}

\begin{document}

\section{Preregistered hypotheses}\label{sec:hypotheses}

\begin{frame}[fragile,allowframebreaks]{Preregistered hypotheses}
The registration in \breaktt{data/example\_registration.json} declares
its hypotheses, outcomes, exclusion rules, and analysis plan before any
result is observed. It is frozen under a SHA-256 content hash (prefix
\texttt{96a34a11d132}); the validator in
\breaktt{registered\_report.protocol.validate\_registration} confirms the
payload carries every required section with no structural findings.
\end{frame}

\begin{frame}[fragile,allowframebreaks]{Confirmatory hypothesis}
\protect\phantomsection\label{confirmatory-hypothesis}
\begin{itemize}
\tightlist
\item
  \textbf{H1} --- \emph{The proposed condition improves the primary
  score.} This is the sole confirmatory hypothesis. It maps to the
  registered outcome \texttt{primary\_score}, measured as a mean score
  difference and analysed with a two-sided label-permutation test at
  \texttt{alpha\ =\ 0.05}.
\end{itemize}

Each registered hypothesis is bound to exactly one registered outcome
and one registered analysis. \breakseq{[@fig:hypothesis\_map]} renders that
binding directly from the frozen registration via
\breaktt{hypothesis\_outcome\_map}, so the figure cannot silently drift
from the locked plan.

\begin{figure}
\centering
\pandocbounded{\includegraphics[keepaspectratio,alt={Registered hypothesis-to-outcome-to-analysis mapping for the demonstration registration. Each row is generated by hypothesis\_outcome\_map(registration) in src/registered\_report/demo\_study.py: hypothesis H1 links to the primary\_score outcome, the mean score difference measure, and the two-sided permutation test analysis. Only registered hypotheses appear; exploratory endpoints are excluded by construction.}]{../figures/hypothesis_map.png}}
\caption{Registered hypothesis-to-outcome-to-analysis mapping for the
demonstration registration. Each row is generated by
\protect\breaktt{hypothesis\_outcome\_map(registration)} in
\protect\breaktt{src/registered\_report/demo\_study.py}: hypothesis H1
links to the \protect\texttt{primary\_score} outcome, the
\protect\breaktt{mean\ score\ difference} measure, and the
\protect\breaktt{two-sided\ permutation\ test} analysis. Only registered
hypotheses appear; exploratory endpoints are excluded by
construction.}\label{fig:hypothesis_map}
\end{figure}
\end{frame}

\begin{frame}[fragile,allowframebreaks]{Registered outcome and exclusion
rule}
\protect\phantomsection\label{registered-outcome-and-exclusion-rule}
\begin{itemize}
\tightlist
\item
  \textbf{Primary outcome} \texttt{primary\_score} --- mean score
  difference between the proposed and reference conditions, evaluated by
  a two-sided permutation test.
\item
  \textbf{Exclusion rule} --- records missing the primary outcome are
  excluded \emph{before} group assignment is inspected, so exclusion
  decisions cannot depend on the observed effect.
\end{itemize}

No secondary outcome is preregistered. Any additional endpoint
introduced during analysis is therefore exploratory and must be recorded
in the deviation ledger ({[}Section @sec:deviations{]}) rather than
presented as a confirmatory finding.
\end{frame}

\end{document}
